% [ILLUSTRATION À GÉNÉRER : aquarelle d'une table d'autopsie ou d'un phare éteint — à créer via Midjourney puis insérer avec \insertImage]

\chapter{Autopsies : trois retours en arrière célèbres}

\lettrine[lines=2]{J}{usqu'ici}, nous avons parlé des échecs en général. Il est temps de nommer des noms — pas pour le plaisir de la démolition, mais parce que trois entreprises célèbres ont eu l'élégance rare de reculer \emph{en public}, en laissant des traces documentées. Ces traces valent de l'or : ce sont les seules boîtes noires dont nous disposons.

Trois crashs, trois causes différentes. Prenez des notes.

\section{GitHub : le flou qui profite aux plus forts}

On a croisé GitHub au début de ce livre ; il faut maintenant finir l'histoire, parce que sa leçon est la plus dérangeante des trois.

GitHub, au début des années 2010, c'est l'utopie du sans-chef version développeurs : pas de managers, chacun choisit ses projets, la légitimité vient du code. Puis, en 2014, une crise interne retentissante éclate — des accusations publiques de harcèlement portées par une développeuse, et l'impossibilité, précisément, de savoir qui aurait dû traiter le problème. Le co-fondateur emblématique démissionne, et l'entreprise réintroduit des managers dans la foulée\footnote{Foss, N. J., \& Klein, P. G. (2014). Why Managers Still Matter. MIT Sloan Management Review, 56(1), 73-80.}.

La leçon n'est pas « le sans-chef mène au harcèlement ». La leçon, c'est que l'absence de structure formelle n'avait pas supprimé le pouvoir chez GitHub — elle l'avait rendu invisible. Des cercles informels décidaient de qui comptait et de qui ne comptait pas, sans qu'aucun organigramme ne permette de les contester, et sans qu'aucun circuit ne protège les plus vulnérables. Quand le pire est arrivé, il n'y avait ni responsable désigné, ni procédure, ni recours. On retrouvera cet argument au chapitre des critiques ; GitHub en est la démonstration grandeur nature.

\section{Zappos : la liberté au forceps}

Zappos, le e-commerçant américain à la culture d'entreprise légendaire, est le plus grand laboratoire d'holacratie de l'histoire — et son récit tient en un paradoxe : on n'impose pas l'auto-organisation par ultimatum.

En 2013-2014, Tony Hsieh, son dirigeant charismatique, décide de basculer l'entreprise entière — environ 1 500 personnes — en holacratie. La transition patine, la confusion s'installe. En mars 2015, Hsieh envoie alors un mémo resté célèbre, « The Offer » : adhérez pleinement à l'auto-organisation, ou partez avec une indemnité. Environ 14 \% des salariés — plus de 200 personnes — prennent l'offre immédiatement ; le décompte grimpera à près de 18 \% des effectifs dans l'année\footnote{Bernstein, E., Bunch, J., Canner, N., \& Lee, M. (2016). Beyond the Holacracy Hype. Harvard Business Review, 94(7-8), 38-49.}.

Arrêtons-nous sur ce chiffre, parce qu'il est le plus honnête de toute la littérature sur la libération : quand on donne réellement le choix, avec un chèque à la clé, presque un salarié sur cinq préfère partir. Voilà la proportion de gens que l'enthousiasme des conférences rend invisibles. Et notons l'ironie de la méthode : le passage à l'auto-organisation décrété d'en haut, par ultimatum du patron — c'est-à-dire l'exact geste managérial que le modèle prétend abolir. Zappos a continué l'expérience des années durant, en l'adaptant considérablement ; mais le mémo de 2015 reste le rappel le plus brutal de ce que ce livre répète : la liberté imposée n'est pas de la liberté.

\section{Medium : mort par excès de structure}

Le troisième cas est mon préféré, parce qu'il échoue par le côté opposé — et qu'il a été documenté par l'entreprise elle-même, avec une transparence exemplaire.

Medium, la plateforme d'écriture fondée par Evan Williams (co-fondateur de Twitter), adopte l'holacratie à ses débuts, en pionnière convaincue. En mars 2016, son directeur des opérations, Andy Doyle, publie un billet expliquant pourquoi l'entreprise arrête : la coordination à l'échelle devenait laborieuse, la tenue des registres de gouvernance chronophage, et surtout — c'est la phrase clé — la codification détaillée des responsabilités finissait par \emph{freiner l'esprit d'initiative et le sentiment d'appropriation collective} que le système était censé produire. L'holacratie, écrit-il en substance, « se mettait en travers du travail »\footnote{Doyle, A. (2016). Management and Organization at Medium. The Medium Blog, mars 2016.}.

Relisez le paradoxe de GitHub, puis celui de Medium : le premier meurt de trop peu de structure, le second de trop. Ce n'est pas une contradiction — c'est la cartographie du chemin de crête. La structure est un médicament : sous-dosée, la maladie du pouvoir informel prospère ; surdosée, elle tue l'élan qu'elle devait protéger. Et le bon dosage n'est écrit dans aucune constitution : il se trouve par l'expérimentation, entreprise par entreprise, année après année.

\section{Ce que les trois autopsies disent ensemble}

Trois entreprises brillantes, riches, pleines de talents — trois retours en arrière. Si vous cherchiez une preuve que la libération ne pardonne pas l'à-peu-près, la voilà. Mais regardez mieux : aucune des trois n'est revenue au taylorisme. GitHub a gardé une culture d'autonomie forte sous ses managers ; Zappos a poursuivi et adapté son expérience pendant des années ; Medium a explicitement conservé les acquis du système en abandonnant sa mécanique\footnote{Bernstein, E., Bunch, J., Canner, N., \& Lee, M. (2016). Beyond the Holacracy Hype. Harvard Business Review, 94(7-8), 38-49.}.

Autrement dit : même les échecs les plus retentissants du mouvement ne sont pas des retours à la case départ. Ce sont des réglages de dosage. C'est, au fond, une nouvelle assez encourageante — à condition d'accepter l'idée que votre libération, comme les leurs, ne ressemblera jamais tout à fait au livre qui vous l'a fait rêver.
